% !TEX root = sommaire.tex

%%------------------------------------------------ Le résumé et les mots-clés -
\begin{ResumeMotsCles}
%% Résumé en français : pas plus de 2000 caractères
\begin{resume}
La résolution de classes de types est un composant central de l'élaboration dans
les assistants de preuve modernes. Cette thèse étudie un moteur de résolution
pour \rocq fondé sur \elpi, un interprète embarqué de \lprolog. Nous
présentons un schéma de compilation traduisant les classes de types et leurs
instances en une base de données de programmation logique, permettant de
déléguer la résolution à \elpi, et évaluons cette approche par rapport au
solveur natif de \rocq.

Une difficulté majeure provient des différences entre les mécanismes
d'unification de \rocq et d'\elpi. Malgré leur prise en charge de l'unification
d'ordre supérieur et de la conversion $\beta\eta$, leurs représentations
internes peuvent entraîner des échecs d'unification après traduction. Pour y
remédier, nous introduisons une technique de compilation fondée sur des
contraintes différées, vérifiées à l'exécution.

La thèse présente également une analyse statique du déterminisme pour \elpi,
capable de vérifier qu'un prédicat produit au plus une solution. Afin d'en
garantir la correction, nous en formalisons les propriétés essentielles dans
\rocq à l'aide d'un modèle du premier ordre de l'interprète \elpi.

Ensemble, ces contributions montrent comment les techniques de programmation
logique peuvent servir à implémenter, analyser et vérifier des mécanismes
d'élaboration dans les assistants de preuve.
\end{resume}

%% Résumé en anglais : pas plus de 2000 caractères
\begin{abstract}

Type-class resolution is a central component of elaboration in modern proof
assistants. This thesis investigates a type-class resolution engine for the \rocq
proof assistant based on \elpi, an embeddable \lprolog interpreter. We
present a compilation scheme translating \rocq type classes and instances into a
logic-programming database, enabling resolution to be delegated to \elpi, and
evaluate its expressiveness and performance against \rocq's native solver.

A key challenge arises from differences between \rocq and \elpi unification.
Although both support higher-order unification and $\beta\eta$-conversion,
differences in term representations can lead to unification failures after
translation. To address this issue, we introduce a compilation technique based
on delayed constraints, replacing problematic subterms with links that record
constraints checked during execution.

The thesis also presents a static determinacy analysis for \elpi, capable of
verifying that selected predicates produce at most one solution. To increase
confidence in the analysis, we formalize its core correctness properties in \rocq
by mechanizing a first-order model of the \elpi interpreter and proving the
soundness of the checker.

Together, these contributions demonstrate how logic-programming techniques can
be used to implement, analyze, and verify elaboration mechanisms in proof
assistants.
\end{abstract}

%% Mots-clés en français
\motscles{Programmation logique, Théorie des types, Elpi, Rocq, Classes de type, Analyse de déterminisme.}
%% Mots-clés en anglais
\keywords{Logic Programming, Type Theory, Elpi, Rocq, Type Classes, Determinacy Analysis.}
\end{ResumeMotsCles}


% Changer la langue du document en anglais:
% - commenter les "renewcommand" dans [ce fichier](./francisation.sty)
% - changer l'ordre des argument du pacquet `babel`: `\RequirePackage[french,english]{babel}` dans [ce fichier](./these-ISSS.cls)
% - changer l'option `selectlanguage`: `\selectlanguage{english}` dans [ce fichier](./these-ISSS.cls)
